\section{Numeral calculations as annotated arithmetic proofs}
\label{sec:p05-pa-numeral-proof-builders}

The following constructions produce actual row trees for the selected
\(\mathsf{Prf}^{19}_{\PA}\) of
\cref{def:p05-selected-direct-pa-graph}. Their verification takes place in
\(\PA\), using the specified graph equations. The output includes the
universal-elimination rows of every arithmetic axiom instance. In particular,
the construction does not enlarge the seven-axiom leaf convention of
\cref{sec:p05-standard-pa-checker}.

\begin{hbdefinition}{def:p05-numeral-builder-notation}{Code operations and fixed descriptions}
Use the descriptions and expansion order of
\cref{def:p05-checker-description-language,def:p05-checker-list-description}.
The following names abbreviate constructor expressions:
\[
\begin{aligned}
 v_i&\defeq N_0(i),& o&\defeq N_1(),& s(t)&\defeq N_2(t),\\
 a(t,u)&\defeq N_3(t,u),& b(t,u)&\defeq N_4(t,u),&
 \mathsf{Eq}(t,u)&\defeq N_5(t,u),\\
 \mathsf{bc}&\defeq N_6(),& \mathsf{Neg}(c)&\defeq N_7(c),&
 c\mathbin{\&}d&\defeq N_8(c,d),\\
 c\Rightarrow d&\defeq N_{10}(c,d),&
 \mathsf{All}(i,c)&\defeq N_{11}(i,c),&
 \mathsf{Ex}(i,c)&\defeq N_{12}(i,c).
\end{aligned}
\]
Here \(\mathsf{bc}\) is only a local code name, not an extra logical constant.
The displayed operations construct syntax codes. Numerical addition and
multiplication retain the names \(A,M\). Write
\(\mathsf{entry}(R,j)\defeq h(T(R,j))\) for the existing list-entry
description, and \(\mathsf{nodes}(c)\) for the already fixed syntax-node-count
description. A reference to an earlier library function inserts its previously
fixed description with its original arguments.

Define, by primitive recursion with its counter last,
\[
 \mathsf{it}(t,0)=t,\qquad
 \mathsf{it}(t,Sj)=s(\mathsf{it}(t,j)),\qquad
 \nu(j)\defeq\mathsf{it}(o,j).
\]
Write \(\mathsf{pack}(R)\defeq N_{19}(N_{15}(R),0)\).
An application of a named description inside arithmetic always means its
unique \(\Gamma\)-graph output, under
\cref{thm:p05-pa-pr-graphs}. Finite displayed templates expand by composition;
every recursion, fold, branch order and failure value is specified below.
Numerical guards \(x=y,x\ne y,x<y,x\leq y\) expand respectively to
\(E(x,y),Z(E(x,y)),\mathsf{lt}(x,y),Z(\mathsf{lt}(y,x))\);
reverse the arguments for \(>\) and \(\geq\). Thus every guarded
success/failure case below is a composition with the already fixed numerical
conditional, including its values on all failing inputs.
\end{hbdefinition}

\begin{hbdefinition}{def:p05-builder-nocapture}{A table for substitution without renaming}
For parameters \(i,u\), form a table with rows \([d,z,w]\), in increasing
expression-code order. A row at \(c\) is zero in all three fields if \(c\)
is neither a term nor a formula, or \(u\) is not a term. Otherwise use the
following ordered constructor cases. Child fields are entries in the earlier
table, exactly as in the selected syntax tables.
\begin{enumerate}
\item A variable \(v_j\) gives \([1,I(E(j,i),u,c),1]\).
Zero and bottom give \([1,c,1]\).
\item A unary nonbinding constructor \(N_k(c_0)\), for \(k=2,7\),
gives \([d_0,N_k(z_0),2+w_0]\). A binary nonbinding constructor
\(N_k(c_0,c_1)\), for \(k=3,4,5,8,9,10\), gives
\([d_0d_1,N_k(z_0,z_1),2+w_0+w_1]\).
\item A binder \(N_k(j,c_0)\), for \(k=11,12\), gives
\([1,c,1]\) when \(j=i\), and otherwise gives
\([d_0Z(\mathsf{AV}(u,j)),N_k(j,z_0),2+w_0]\).
\end{enumerate}
In each case replace the whole row by \([0,0,0]\) if its first field is
zero. Retrieve row \(c\) after \(Sc\) updates and call its fields
\(\mathsf{nc}(c,i,u),\mathsf{ns}(c,i,u),\mathsf{cost}(c,i,u)\).
These are three projections of one fixed finite table description. Wrong
tags, arities or child types are already excluded by its initial type test.
\end{hbdefinition}

\begin{hblemma}{thm:p05-builder-nocapture}{proved=annex:p5.builder-nocapture}{Internal substitution identities needed by the builders}
In \(\PA\), if \(\mathsf{nc}(c,i,u)=1\), the table output has the
input expression's term or formula type and
\[
 \mathsf{sub}(c,i,u)=\mathsf{ns}(c,i,u),\qquad
 1\leq\mathsf{cost}(c,i,u)\leq2\mathsf{nodes}(c).
\]
Thus its table equations are internal equations for the selected substitution
machine on these inputs. If \(u\) is a term with no variable occurrences,
the first field is one for every valid \(c\). It is also one for every
quantifier-free valid expression and every term \(u\).

In particular \(\PA\) proves that \(\nu(j)\) is a variable-free term
and, assuming \(\mathsf{Tm}(t)=1\), that
\(\mathsf{it}(s(t),j)=\mathsf{it}(t,Sj)\). With universal numerical
closure it also proves the guarded numeral-invariance assertion
\[
 \mathsf{Tm}(u)=1\ \longrightarrow
 \bigl(\mathsf{ns}(\nu(n),i,u)=\nu(n)\land
        \mathsf{sub}(\nu(n),i,u)=\nu(n)\bigr).
\]
The finite open-term axiom instances used below satisfy the same
table test; their only surviving inner binder has index 1 and their outer
replacement contains no variable of that index.
\end{hblemma}

\begin{hbdefinition}{def:p05-builder-row-templates}{Exact row joins and elementary proof templates}
All builders first construct raw row lists. For a fixed arity
\(0\leq h\leq3\), let \(R_0,\ldots,R_{h-1}\) be their input lists,
\(l_j=\ell(R_j)\), and \(t_j=\sum_{k<j}l_k\). Define
\(\mathsf{J}_r(c,U;R_0,\ldots,R_{h-1})\) by concatenating, in that
order, \(\mathsf{SR}(t_j,R_j)\), then appending the single row
\[
 [r,c,[t_0+D(l_0),\ldots,t_{h-1}+D(l_{h-1})],U,0].
\]
The empty concatenation and empty reference list are both zero. This is a
separate fixed composition for each displayed arity. It is total on raw
lists; validity claims require every supplied premise list to be nonempty.
Define
\[
 \mathsf{L}(c,l)\defeq[[0,c,0,0,l]],\qquad
 \mathsf{R}(t)\defeq\mathsf{J}_{18}(\mathsf{Eq}(t,t),[t];\ ),
 \qquad f(R)\defeq\mathsf{root}(R).
\]
The second argument of \(\mathsf L\) is the stored marker: marker \(Si\)
means label \(i\), whereas marker zero is unmarked.

The following table specifies all remaining primitive templates. Symbols
\(c,d,t,u,i,y,k\) in its parameter columns are supplied numerical inputs;
each formula/term expression means the constructor expression already fixed.
\[
\begin{array}{c|c}
\text{template}&\text{row-list expression}\\\hline
 \mathsf{FE}(R,t)&
 \mathsf J_{13}(\mathsf{sub}(c,i,t),[i,c,t];R)
       \quad(f(R)=\mathsf{All}(i,c))\\
 \mathsf{FI}(R,i)&\mathsf J_{15}(\mathsf{All}(i,f(R)),[i];R)\\
 \mathsf{EI}(R,i,c,t)&\mathsf J_{14}(\mathsf{Ex}(i,c),[i,c,t];R)\\
 \mathsf{EE}(P,R,i,c,y,k)&\mathsf J_{16}(f(R),[i,c,y,k];P,R)\\
 \mathsf{II}(R,c,k)&\mathsf J_7(c\Rightarrow f(R),[k];R)\\
 \mathsf{NI}(R,c,k)&\mathsf J_9(\mathsf{Neg}(c),[k];R)\\
 \mathsf{B}(R,S)&\mathsf J_{10}(\mathsf{bc},0;R,S)\\
 \mathsf{And}(R,S)&\mathsf J_1(f(R)\mathbin{\&}f(S),0;R,S)\\
 \mathsf{MP}^{\rm row}(R,S)&\mathsf J_8(d,0;R,S)
                         \quad(f(S)=c\Rightarrow d)\\
 \mathsf{EEq}(R,S,t,u,c,i)&
       \mathsf J_{19}(\mathsf{sub}(c,i,u),[t,u,c,i];R,S).
\end{array}
\]
The two patterns in this table are decoded by the actual tag/arity
projections; their template returns zero if that pattern fails. All other
templates have the displayed value on every numerical input. Their invariants
include precisely the appropriate rule premises and side conditions.

If \(R\) has root \(\mathsf{Eq}(t,u)\), let \(j=1+\max(m(t),m(u))\) and set
\[
 \mathsf{Sym}(R)\defeq
 \mathsf{EEq}(R,\mathsf R(t),t,u,\mathsf{Eq}(v_j,t),j).
\]
If \(R,S\) have roots \(\mathsf{Eq}(t,u),\mathsf{Eq}(u,w)\), let
\(j=1+\max(m(t),m(u),m(w))\) and set
\[
 \mathsf{Tr}(R,S)\defeq\mathsf{EEq}(S,R,u,w,\mathsf{Eq}(t,v_j),j).
\]
For symmetry, decode both endpoints from its sole equality root. For
transitivity, take \(t,u\) from the first and second arguments of \(f(R)\)
and \(w\) from the second argument of \(f(S)\). Test that each indicated
root has tag 5 and arity two; failure returns zero. Equality of the first
argument of \(f(S)\) with \(u\) is an invariant hypothesis, not an
additional branch test in this description.

For each fixed one-hole term context \(g\) used below, define
\(\mathsf{Cg}_g(R)\), for root \(\mathsf{Eq}(t,u)\), by
\[
 \mathsf{EEq}(R,\mathsf R(g(t)),t,u,\mathsf{Eq}(g(t),g(v_j)),j),
 \quad j=1+\max(m(t),m(u),m(g(t)),m(g(u))).
\]
Here the contexts are exactly successor, left addition by a supplied term,
right addition by a supplied term, and the analogous two multiplication
contexts. Each is a separately fixed constructor expression with its fixed
terms as additional parameters. There is no recursion over contexts in this
definition. It returns zero on a failed equality-root pattern.
\end{hbdefinition}

\begin{hblemma}{thm:p05-builder-row-templates}{proved=annex:p5.builder-row-templates}{Internal preservation for raw proof joins}
\(\PA\) proves the following assertions for the fixed templates. Joining
nonempty trees with successful rows, and with the exact local test at the
new row satisfied, gives another tree with all rows successful. Each old
nonroot keeps parent count one; each old root receives one parent; the new
root has count zero. The new open list is exactly the selected row rule's
update on the shifted premise open lists.

Consequently the templates for symmetry, transitivity and congruence have
their displayed equality conclusions. They retain only the old open leaves.
Every retained leaf has the same formula and marker as its original copy.
The quantifier and discharge templates additionally require their actual
free-variable or all-variable tests and a surviving matching marked leaf
where the rule requires nonvacuity. They assert no vacuous discharge rule.
For sentence proofs with axiom leaves, the raw implication-elimination join
is precisely the successful value of \(\mathsf{mp}_{19}\) after packing.
\end{hblemma}

\begin{hbdefinition}{def:p05-builder-axiom-instances}{Axiom instances and two explicit derived constants}
For the seven literal constants \(q_1,\ldots,q_7\), put
\[
 \mathsf A_j(t_0,\ldots,t_{r-1})\defeq
  \mathsf{FE}(\cdots\mathsf{FE}(\mathsf L(q_j,0),t_0)\cdots,t_{r-1}),
\]
where \(r=1\) for \(j=1,3,4,6\), and \(r=2\) for \(j=2,5,7\).
The order is the printed outside-in order, indices 0 then 1 when there are
two binders. No \(q_j\) instance is an assumption leaf.

The next finite expressions construct two derived sentence proofs. Set
\[
\begin{aligned}
 m_Z&\defeq \mathsf{Eq}(a(o,v_0),v_0),& I_Z&\defeq\mathsf{IndAx}(m_Z,0),\\
 m_L&\defeq \mathsf{Eq}(a(s(v_0),v_1),s(a(v_0,v_1))),&
 I_L&\defeq\mathsf{IndAx}(m_L,1).
\end{aligned}
\]
Their concrete matrices and indices are fixed here; their axiom codes use
the existing fully specified canonical-induction constructor. Define
\[
\begin{aligned}
 Z_0&\defeq\mathsf A_4(o),\\
 Z_s&\defeq\mathsf{Tr}(\mathsf A_5(o,v_0),
                         \mathsf{Cg}_s(\mathsf L(m_Z,1))),\\
 D_Z&\defeq\mathsf{MP}^{\rm row}
  (\mathsf{And}(Z_0,\mathsf{FI}(\mathsf{II}(Z_s,m_Z,0),0)),
                                                     \mathsf L(I_Z,0)),\\
 L_0&\defeq\mathsf{Tr}(\mathsf A_4(s(v_0)),
                              \mathsf{Sym}(\mathsf{Cg}_s(\mathsf A_4(v_0)))),\\
 L_s&\defeq\mathsf{Tr}\bigl(
       \mathsf{Tr}(\mathsf A_5(s(v_0),v_1),
                           \mathsf{Cg}_s(\mathsf L(m_L,1))),
                         \mathsf{Sym}(\mathsf{Cg}_s(\mathsf A_5(v_0,v_1)))\bigr),\\
 D_L&\defeq\mathsf{FI}\bigl(
   \mathsf{MP}^{\rm row}
    (\mathsf{And}(L_0,\mathsf{FI}(\mathsf{II}(L_s,m_L,0),1)),
                                   \mathsf{FE}(\mathsf L(I_L,0),v_0)),0\bigr).
\end{aligned}
\]
Each repeated template occurrence inserts a fresh copy of its entire row
tree. Set \(\mathsf A_Z(t)=\mathsf{FE}(D_Z,t)\) and
\(\mathsf A_L(t,u)=\mathsf{FE}(\mathsf{FE}(D_L,t),u)\).
These expressions copy the derived proofs; they do not create new axiom
leaves with their conclusions.
\end{hbdefinition}

\begin{hblemma}{thm:p05-builder-left-addition}{proved=annex:p5.builder-left-addition}{The two derived addition proofs}
The packed records \(D_Z,D_L\) satisfy the selected checker, with roots
\[
 \mathsf{All}(0,m_Z),\qquad
 \mathsf{All}(0,\mathsf{All}(1,m_L)),
\]
respectively. Their decoded conclusions are
\(\forall x_0(0+x_0=x_0)\) and
\(\forall x_0\forall x_1(Sx_0+x_1=S(x_0+x_1))\).
Their surviving leaves are unmarked members of the fixed list
\[
 \mathcal A_*\defeq[q_1,q_2,q_3,q_4,q_5,q_6,q_7,I_Z,I_L].
\]
Every member passes \(\mathsf{Ax}_{19}\), and every variable occurring in
one of these leaves has index 0 or 1. These assertions, including the
checker assertions, hold as theorems of \(\PA\) for the displayed fixed
constructor expressions.
\end{hblemma}

\begin{hbdefinition}{def:p05-numeral-arithmetic-builders}{Arithmetic and successor-string recurrences}
Define the following raw row-list recurrences in their displayed order.
Numerical output values use the already fixed descriptions \(A,M\).
\[
\begin{aligned}
 R_+(m,0)&=\mathsf A_4(\nu(m)),\\
 R_+(m,Sn)&=\mathsf{Tr}(\mathsf A_5(\nu(m),\nu(n)),
                                      \mathsf{Cg}_s(R_+(m,n))),\\
 R_\times(m,0)&=\mathsf A_6(\nu(m)),\\
 R_\times(m,Sn)&=\mathsf{Tr}\bigl(
   \mathsf{Tr}(\mathsf A_7(\nu(m),\nu(n)),
                 \mathsf{Cg}_{z\mapsto a(z,\nu(m))}(R_\times(m,n))),
                                            R_+(M(m,n),m)\bigr),\\
 R_{\rm right}(t,0)&=\mathsf A_4(t),\\
 R_{\rm right}(t,Sn)&=\mathsf{Tr}(\mathsf A_5(t,\nu(n)),
                                \mathsf{Cg}_s(R_{\rm right}(t,n))),\\
 R_{\rm left}(t,0)&=\mathsf A_Z(t),\\
 R_{\rm left}(t,Sn)&=\mathsf{Tr}(\mathsf A_L(\nu(n),t),
                                 \mathsf{Cg}_s(R_{\rm left}(t,n))).
\end{aligned}
\]
The parameter \(t\) in the last two recurrences will be a numeral,
\(v_y\), or \(s(v_y)\), with \(y\geq3\). These are their asserted
domains. On other inputs the same equations give total numerical outputs,
with the already specified template failure values.
\end{hbdefinition}

\begin{hbtheorem}{thm:p05-numeral-arithmetic-builders}{proved=annex:p5.numeral-arithmetic-builders}{Internal arithmetic-builder invariants}
With universal numerical quantification, \(\PA\) proves
\[
\begin{aligned}
 \mathsf{Prf}^{19}_{\PA}
  (\mathsf{pack}(R_+(m,n)),\mathsf{Eq}(a(\nu(m),\nu(n)),\nu(A(m,n)))),\\
 \mathsf{Prf}^{19}_{\PA}
  (\mathsf{pack}(R_\times(m,n)),\mathsf{Eq}(b(\nu(m),\nu(n)),\nu(M(m,n)))).
\end{aligned}
\]
For every permitted \(t\), the two remaining recurrences produce successful
trees with roots
\[
 \mathsf{Eq}(a(t,\nu(n)),\mathsf{it}(t,n)),\qquad
 \mathsf{Eq}(a(\nu(n),t),\mathsf{it}(t,n)),
\]
respectively. Their surviving leaves, and those of the numeral recurrences,
are unmarked members of \(\mathcal A_*\). A root containing \(v_y\)
is an intermediate open formula; it is not asserted to pass the sentence
test until the later construction closes it.
\end{hbtheorem}

\begin{hbdefinition}{def:p05-numeral-literal-builders}{Disequality and both comparison polarities}
First define a contradiction constructor. For \(t=o\) or \(t=v_y\)
with \(y\geq3\), use numerical parameters \(k,b\) and a raw tree \(R\).
The counter \(j\) is last:
\[
\begin{aligned}
 P(t,k,b,R,0)&=R,\\
 P(t,k,b,R,Sj)&=
 \begin{cases}
 \mathsf{MP}^{\rm row}\bigl(P(t,k,b,R,j),
    \mathsf A_2(\mathsf{it}(t,L(k,Sj)),\nu(L(b,Sj)))\bigr),&j<b,\\
 P(t,k,b,R,j),&j\geq b,
 \end{cases}\\
 \mathsf{Clash}(t,k,b,R)&\defeq
 \mathsf B\bigl(P(t,k,b,R,b),
                    \mathsf A_1(\mathsf{it}(t,D(L(k,b))))\bigr).
\end{aligned}
\]
The two cases expand to the numerical conditional in their printed order.
For the stated domain of \(t\), the asserted success premises are
\[
 \mathsf{Tree}(R)=1,\qquad k>b,\qquad
 f(R)=\mathsf{Eq}(\mathsf{it}(t,k),\nu(b)).
\]

For \(m\ne n\), let \(c=\mathsf{Eq}(\nu(m),\nu(n))\), and set
\[
 R_{\ne}(m,n)\defeq
 \begin{cases}
 \mathsf{NI}(\mathsf{Clash}(o,m,n,\mathsf L(c,1)),c,0),&n<m,\\
 \mathsf{NI}(\mathsf{Clash}(o,n,m,\mathsf{Sym}(\mathsf L(c,1))),c,0),&m<n,\\
 0,&m=n.
 \end{cases}
\]
Set \(R_{=}(m,n)=I(E(m,n),\mathsf R(\nu(m)),0)\).

The inherited comparison expansions retain a supplied binder index \(i\):
\[
 \mathsf{Le}_i(t,u)\defeq\mathsf{Ex}(i,\mathsf{Eq}(a(v_i,t),u)),\qquad
 \mathsf{Lt}_i(t,u)\defeq\mathsf{Ex}(i,\mathsf{Eq}(a(t,s(v_i)),u)).
\]
For closed numeral terms define, in the success cases,
\[
\begin{aligned}
 R_{\leq}^+(i,m,n)&\defeq
  \mathsf{EI}(R_+(L(n,m),m),i,\mathsf{Eq}(a(v_i,\nu(m)),\nu(n)),\nu(L(n,m))),
       &&m\leq n,\\
 R_<^+(i,m,n)&\defeq
  \mathsf{EI}(R_+(m,L(n,m)),i,\mathsf{Eq}(a(\nu(m),s(v_i)),\nu(n)),
                                            \nu(D(L(n,m)))),&&m<n.
\end{aligned}
\]
Outside the indicated numerical condition each builder returns zero.

For either negative comparison put \(y=i+3\), let \(c\) be the matrix
of that comparison at \(\nu(m),\nu(n)\), and let
\(C=\mathsf{Ex}(i,c)\). Take
\[
 H=\mathsf L(C,1),\qquad
 K=\mathsf L(\mathsf{sub}(c,i,v_y),2).
\]
For \(m>n\), set
\[
 V_{\leq}=\mathsf{Tr}(\mathsf{Sym}(R_{\rm right}(v_y,m)),K),\qquad
 W_{\leq}=\mathsf{Clash}(v_y,m,n,V_{\leq}).
\]
For \(m\geq n\), set
\[
 V_<=\mathsf{Tr}(\mathsf{Sym}(R_{\rm left}(s(v_y),m)),K),\qquad
 W_<=\mathsf{Clash}(v_y,Sm,n,V_<).
\]
In the respective success case define
\[
 R_\diamond^-(i,m,n)\defeq
 \mathsf{NI}(\mathsf{EE}(H,W_\diamond,i,c,y,1),C,0),
                 \qquad\diamond\in\{\leq,<\}.
\]
The failure value is zero. Thus the existential assumption has label 0,
its minor instance has label 1, and the eigenvariable is exactly \(i+3\).
Finally put
\[
 R_{\neg\bot}\defeq\mathsf{NI}(\mathsf L(\mathsf{bc},1),\mathsf{bc},0).
\]
\end{hbdefinition}

\begin{hbtheorem}{thm:p05-numeral-literal-builders}{proved=annex:p5.numeral-literal-builders}{Internal literal-proof assertions}
\(\PA\) proves, with all numerical parameters universally quantified,
that the packed builders in the following table satisfy
\(\mathsf{Prf}^{19}_{\PA}\) with exactly the listed root whenever its
condition holds:
\[
\begin{array}{c|c|c}
 \text{raw builder}&\text{condition}&\text{literal root}\\\hline
 R_=(m,n)&m=n&\mathsf{Eq}(\nu(m),\nu(n))\\
 R_{\ne}(m,n)&m\ne n&\mathsf{Neg}(\mathsf{Eq}(\nu(m),\nu(n)))\\
 R_{\leq}^+(i,m,n)&m\leq n&\mathsf{Le}_i(\nu(m),\nu(n))\\
 R_{\leq}^-(i,m,n)&m>n&\mathsf{Neg}(\mathsf{Le}_i(\nu(m),\nu(n)))\\
 R_<^+(i,m,n)&m<n&\mathsf{Lt}_i(\nu(m),\nu(n))\\
 R_<^-(i,m,n)&m\geq n&\mathsf{Neg}(\mathsf{Lt}_i(\nu(m),\nu(n)))\\
 R_{\neg\bot}&\text{always}&\mathsf{Neg}(\mathsf{bc}).
\end{array}
\]
All surviving leaves are unmarked members of \(\mathcal A_*\).
Under the stated term domain and the premises \(\mathsf{Tree}(R)=1\),
\(k>b\) and \(f(R)=\mathsf{Eq}(\mathsf{it}(t,k),\nu(b))\), the
contradiction constructor has successful rows, root \(\mathsf{bc}\), and
retains its input open leaves together with unmarked \(q_1,q_2\) leaves.
These are internal checker invariants, including
zero numeral arguments, a zero peeling bound, equal comparison endpoints,
and the premise-free identity case.
\end{hbtheorem}

\begin{hbdefinition}{def:p05-fixed-term-literal-builders}{A builder for each fixed term and literal}
Fix a finite list \([i_0,\ldots,i_{r-1}]\) of distinct variable indices
and a term \(t\) whose variables belong to that list. The variable
\(x_{i_j}\) receives the numerical input \(m_j\). For this numerical tuple
\(\bar m\), define its value \(t^\#(\bar m)\) by the fixed term
formation using numerical \(0,S,A,M\). Define its numeral-instance code
\(t^\nu(\bar m)\) by the same formation using \(o,s,a,b\) and the
numeral codes of the supplied entries.

Define a raw proof \(R_t(\bar m)\) by meta-recursion on this one fixed
term. Zero and a variable use \(\mathsf R(o)\) and
\(\mathsf R(\nu(m_j))\), respectively. Successor uses
\(\mathsf{Cg}_s(R_u)\). For a sum of fixed terms \(u,v\), use
\[
 \mathsf{Tr}\bigl(
  \mathsf{Tr}(\mathsf{Cg}_{z\mapsto a(z,v^\nu)}(R_u),
       \mathsf{Cg}_{z\mapsto a(\nu(u^\#),z)}(R_v)),
                              R_+(u^\#,v^\#)\bigr).
\]
For a product replace both addition contexts by the corresponding
multiplication contexts and the final builder by \(R_\times\).
All tuple arguments in this display are unchanged \(\bar m\).
This assigns a separately fixed finite description to each term; the
object-language parameters never range over term descriptions.

For a fixed signed equality or signed comparison of terms \(u,v\), write
\(W(t,z)\) for its formula-code constructor, including its sign. In a
comparison retain the binder \(i\) selected when the original abbreviation
was expanded: do not select a new binder after numeral substitution.
For equality use the auxiliary numerical value \(i=0\). Choose
\(j=1+\max(i,m(u^\nu),m(v^\nu))\), and take the appropriate normalized
raw literal proof \(Q\) at \(u^\#,v^\#\) from the preceding table.
Apply successively
\[
\begin{aligned}
 Q_1&=\mathsf{EEq}(\mathsf{Sym}(R_u),Q,
                   \nu(u^\#),u^\nu,W(v_j,\nu(v^\#)),j),\\
 Q_2&=\mathsf{EEq}(\mathsf{Sym}(R_v),Q_1,
                   \nu(v^\#),v^\nu,W(u^\nu,v_j),j).
\end{aligned}
\]
The final description returns \(\mathsf{pack}(Q_2)\) when the numerical
condition for that signed literal holds, and zero otherwise. The literal
\(\neg\bottom\) uses \(\mathsf{pack}(R_{\neg\bot})\); the positive
bottom builder returns zero on every input.
\end{hbdefinition}

\begin{hbtheorem}{thm:p05-fixed-term-literal-builders}{proved=annex:p5.fixed-term-literal-builders}{Fixed-term evaluation and literal transport inside PA}
For each fixed term \(t\), \(\PA\) proves that
\(\mathsf{pack}(R_t(\bar m))\) satisfies the selected checker with root
\[
                 \mathsf{Eq}(t^\nu(\bar m),\nu(t^\#(\bar m))).
\]
For each fixed literal \(\lambda\) just specified, with final builder
description \(d_\lambda\) and literal numeral-instance code
\(\lambda^\nu(\bar m)\), it proves
\[
 \forall\bar m\,\bigl(\lambda(\bar m)\to
   \exists p\,(\Gamma_{d_\lambda}(\bar m,p)\land
                \mathsf{Prf}^{19}_{\PA}(p,\lambda^\nu(\bar m)))\bigr).
\]
In this arithmetic implication \(\lambda(\bar m)\) denotes the original
fixed formula at the numerical parameters, not its code or an internal truth
predicate. All code-valued applications are eliminated through their specified
graphs. The positive-bottom instance has contradictory antecedent.
\end{hbtheorem}

\paragraph{The next construction.}
The next component combines these literal builders through the connectives and
marked bounded quantifiers of each fixed bounded formula. It must still
construct finite case enumeration, convert it to a quantified derivation, and
verify dependent substitutions and all discharge/freshness conditions. The
literal construction above supplies its arithmetic leaves. It does not itself
establish the bounded-formula compiler, the evidence-to-proof theorem, D3 or
the remaining provability-logic results.
